% ============================================================================
% DRAFT (2026-07-18) — Νέο κεφάλαιο αποτελεσμάτων: leak-free campaign
% ΔΕΝ έχει μπει στο main.tex. Αντικαθιστά (μετά από έγκριση) το παλιό
% experiments.tex που αφορά το προ-TZFIX/προ-AEL campaign Δεκ-2025.
% Κάθε αριθμός εδώ έχει trace: βλ. σχόλια %TRACE δίπλα σε κάθε πίνακα.
% ============================================================================
\chapter{Αποτελέσματα Αξιολόγησης — Leak-free Campaign}
\label{ch:experiments_lf}

Το κεφάλαιο αυτό παρουσιάζει τα αποτελέσματα της τελικής πειραματικής εκστρατείας,
η οποία διεξήχθη μετά από δύο δομικές διορθώσεις εγκυρότητας στο pipeline: την
ευθυγράμμιση των τεσσάρων χρονικών ``ρολογιών'' του συνόλου δεδομένων (ενότητα
\ref{sec:lf_validity}) και την επιβολή αυστηρού ελέγχου διαθεσιμότητας πληροφορίας
στα αναδρομικά χαρακτηριστικά (Availability Enforcement Layer, AEL). Τα
αποτελέσματα προγενέστερων εκδόσεων του pipeline --- συμπεριλαμβανομένων όσων
βασίζονταν σε same-day διασυνοριακές τιμές ή σε μη ευθυγραμμισμένα δεδομένα ---
έχουν αποσυρθεί και δεν συγκρίνονται με τα παρόντα.

\section{Πρωτόκολλο Εγκυρότητας}
\label{sec:lf_validity}

Η αξιολόγηση διέπεται από ένα ρητό πρωτόκολλο εγκυρότητας (validity gate), το οποίο
προηγείται κάθε ισχυρισμού απόδοσης:

\begin{itemize}
  \item \textbf{Πύλη πληροφορίας (strict gate).} Κάθε χαρακτηριστικό εισέρχεται στο
        μοντέλο μόνο εφόσον ήταν δημοσιευμένο πριν από το κλείσιμο της αγοράς
        (12:00 CET D-1 για την DAM). Για κάθε νέα πηγή δεδομένων απαιτείται γραπτή
        τεκμηρίωση του χρόνου δημοσίευσης, έλεγχος ετεροσυσχέτισης (lag-scan) και
        έλεγχος ωριαίου προφίλ, πριν από οποιοδήποτε πείραμα.
  \item \textbf{Χρονική ευθυγράμμιση (TZFIX).} Εντοπίστηκε και διορθώθηκε συστημική
        απόκλιση ζωνών ώρας μεταξύ τεσσάρων οικογενειών στηλών του dataset (η
        πραγματική παραγωγή ήταν μετατοπισμένη $+1$h, οι προβλέψεις ΑΠΕ $1$--$2$h
        νωρίς, τα μετεωρολογικά $1$h νωρίς το καλοκαίρι). Όλα τα δεδομένα
        επανακατασκευάστηκαν σε ενιαίο πλαίσιο CET/CEST.
  \item \textbf{AEL (freeze-at-cutoff).} Τα αναδρομικά χαρακτηριστικά πραγματικών
        μεγεθών (φορτίο, παραγωγή, υπολειπόμενο φορτίο) ``παγώνουν'' στην τελευταία
        τιμή που ήταν γνωστή κατά το κλείσιμο της πύλης. Η ορθότητα ελέγχεται με
        δηλητηρίαση δεδομένων (poisoning tests): τεχνητή αλλοίωση τιμών μετά το
        cutoff πρέπει να αφήνει τις προβλέψεις αμετάβλητες (διαφορά $0{,}000000$),
        ενώ αλλοίωση νόμιμων τιμών πρέπει να τις μεταβάλλει. Οι έλεγχοι αυτοί
        εκτελούνται αυτόματα πριν από κάθε πειραματική παρτίδα.
  \item \textbf{Κριτήρια αποδοχής.} Ένα εύρημα γίνεται δεκτό μόνο εφόσον
        $|\Delta \mathrm{MAE}| > 0{,}15$ με συνεπές πρόσημο σε τουλάχιστον δύο
        ανεξάρτητα χρονικά παράθυρα· τα κεντρικά αποτελέσματα απαιτούν επιπλέον
        τρεις τυχαίους σπόρους (seeds) και δεύτερο παράθυρο.
  \item \textbf{Χωρίς βελτιστοποίηση υπερπαραμέτρων.} Σε αντίθεση με το προγενέστερο
        στάδιο της εργασίας, δεν χρησιμοποιείται αυτόματη βελτιστοποίηση
        υπερπαραμέτρων· η ποικιλομορφία επιδιώκεται μέσω των συνόλων χαρακτηριστικών
        και των στρατηγικών, ώστε τα συμπεράσματα να μην εξαρτώνται από tuning
        πάνω στο παράθυρο ελέγχου.
\end{itemize}

Ένα διδακτικό παράδειγμα της αναγκαιότητας του πρωτοκόλλου είναι τα διασυνοριακά
χαρακτηριστικά: οι same-day τιμές γειτονικών ζωνών (IT, BG) εμφάνιζαν εντυπωσιακή
βελτίωση (MAE $14{,}4$--$15{,}0$~€/MWh), η οποία όμως ήταν διαρροή --- οι τιμές
αυτές προκύπτουν από την ίδια δημοπρασία SDAC με τον στόχο και δημοσιεύονται
\emph{μετά} το κλείσιμο της πύλης. Η νόμιμη (καθυστερημένη) εκδοχή τους
αποδείχθηκε επιζήμια τον χειμώνα και απορρίφθηκε από το τελικό σύνολο
χαρακτηριστικών.

\section{Ρύθμιση Πειραμάτων}
\label{sec:lf_setup}

Η αξιολόγηση χρησιμοποιεί walk-forward εβδομαδιαία επανεκπαίδευση (expanding
window) και αναδρομική (recursive) στρατηγική 24ωρης πρόβλεψης, σε δύο ανεξάρτητα
παράθυρα ελέγχου:

\begin{itemize}
  \item \textbf{Q1}: 1 Δεκεμβρίου 2025 -- 28 Φεβρουαρίου 2026 (2.160 ώρες).
  \item \textbf{Καλοκαίρι 2025}: 1 Ιουνίου -- 31 Αυγούστου 2025 (2.208 ώρες).
\end{itemize}

Η εκπαίδευση εκκινεί από τον Ιανουάριο του 2015 και επεκτείνεται έως το εκάστοτε
cutoff κάθε εβδομάδας. Ως μοντέλα αναφοράς χρησιμοποιούνται τα Naive-24,
Naive-168 και το εποχιακό προφίλ τεσσάρων εβδομάδων, υπολογισμένα στα ίδια
παράθυρα. Το Naive-1 δεν αποτελεί έγκυρη αναφορά για ημερήσια πρόβλεψη στο
κλείσιμο της πύλης, καθώς προϋποθέτει πρόσβαση στην πραγματική τιμή της
προηγούμενης ώρας.

%TRACE: υπολογισμός 2026-07-18 από data/processed/hourly.parquet (system python,
%       shift(24)/shift(168)/μέσος 4 εβδομαδιαίων shifts στο ίδιο παράθυρο).
\begin{table}[H]
  \centering
  \caption{Μοντέλα αναφοράς για την τιμή DAM στα δύο leak-free παράθυρα ελέγχου (MAE, €/MWh).}
  \label{tab:lf_benchmarks}
  \begin{tabular}{lcc}
    \toprule
    \textbf{Μοντέλο αναφοράς} & \textbf{Q1 (Δεκ.--Φεβ.)} & \textbf{Καλοκαίρι 2025} \\
    \midrule
    Naive-24            & 21{,}638 & 19{,}111 \\
    Seasonal Profile 4w & 26{,}912 & 19{,}783 \\
    Naive-168           & 27{,}723 & 21{,}975 \\
    \bottomrule
  \end{tabular}
\end{table}

\section{Κεντρικό Αποτέλεσμα: Τιμή DAM}
\label{sec:lf_headline}

Το καλύτερο επικυρωμένο μοντέλο είναι το LightGBM με το σύνολο χαρακτηριστικών
\texttt{default,dense} (ημερολογιακά, υστερήσεις τιμής --- αραιές και πυκνές
1--23h ---, κυλιόμενα στατιστικά, day-ahead προβλέψεις ΑΠΕ/φορτίου, μετεωρολογικά,
καύσιμα), με εβδομαδιαία επανεκπαίδευση και αναδρομική στρατηγική:

%TRACE: runs Block G (2026-07-06): seeds 42/7/123, βλ. ABLATION_PLAN §5.12δ.
%       Q1: 17.035/16.940/16.891 (std 0.060) · Summer: 13.812/13.691/13.708 (std 0.054).
\begin{table}[H]
  \centering
  \caption{Κεντρικό αποτέλεσμα: LightGBM \texttt{default,dense}, εβδομαδιαία
           επανεκπαίδευση, recursive. Μέσο MAE τριών seeds ($\pm$ τυπική απόκλιση).}
  \label{tab:lf_headline}
  \begin{tabular}{lccc}
    \toprule
    \textbf{Παράθυρο} & \textbf{MAE (€/MWh)} & \textbf{std (3 seeds)} & \textbf{Naive-24} \\
    \midrule
    Q1 (Δεκ. 2025--Φεβ. 2026) & 16{,}96 & 0{,}060 & 21{,}638 \\
    Καλοκαίρι 2025            & 13{,}74 & 0{,}054 & 19{,}111 \\
    \bottomrule
  \end{tabular}
\end{table}

Η βελτίωση έναντι του Naive-24 είναι $21{,}6\%$ στο Q1 και $28{,}1\%$ το
καλοκαίρι. Η τυπική απόκλιση μεταξύ seeds ($\leq 0{,}06$~€/MWh) είναι σημαντικά
μικρότερη από τα μεγέθη των υπό εξέταση διαφορών, γεγονός που καθιστά τα
συμπεράσματα των ablation μελετών (ενότητα \ref{sec:lf_ablation}) στατιστικά
διακρίσιμα από τον θόρυβο αρχικοποίησης.

\section{Μελέτες Αφαίρεσης (Ablation) --- Τιμή}
\label{sec:lf_ablation}

Τα κύρια επικυρωμένα ευρήματα (κάθε ένα με συνεπές πρόσημο σε $\geq 2$ ανεξάρτητα
παράθυρα και $\geq 2$ αλγορίθμους όπου σημειώνεται):

\begin{itemize}
  \item \textbf{Πυκνές υστερήσεις (dense, $y_{t-4}\dots y_{t-23}$): θετικές} ---
        4/4 συνθήκες (LGBM+XGB $\times$ Q1+Καλοκαίρι). Συμμετέχουν στο τελικό
        σύνολο χαρακτηριστικών.
  \item \textbf{Μετεωρολογικά: θετικά στη recursive στρατηγική} --- 4/4 συνθήκες.
        Σημειώνεται ρητά ότι τα μετεωρολογικά της τρέχουσας υλοποίησης για την
        τιμή είναι \emph{παρατηρημένες} τιμές (oracle covariate)· η έντιμη
        D-1 εκδοχή (vintage) έχει υλοποιηθεί και επικυρωθεί για το φορτίο
        (ενότητα \ref{sec:lf_load}).
  \item \textbf{Recursive $>$ Direct στην εβδομαδιαία επανεκπαίδευση} --- 4/4
        συνθήκες (Q1+Καλοκαίρι $\times$ LGBM+XGB), διαφορές $-1{,}1$ έως
        $-2{,}9$~€/MWh υπέρ της recursive.
  \item \textbf{Διασυνοριακές τιμές: απορρίφθηκαν} --- same-day εκδοχή
        $=$ διαρροή (ίδια δημοπρασία SDAC)· καθυστερημένη εκδοχή επιζήμια τον
        χειμώνα σε 3/3 ρυθμούς επανεκπαίδευσης.
  \item \textbf{Ensembles: κανένα δεν υπερτερεί του καλύτερου μεμονωμένου
        μοντέλου} υπό έντιμη out-of-sample βαθμονόμηση βαρών.
\end{itemize}

\section{Πρόβλεψη Φορτίου (STLF) --- Σύνοψη Τρέχουσας Εκστρατείας}
\label{sec:lf_load}

% ΣΗΜΕΙΩΣΗ DRAFT: η load εκστρατεία είναι σε εξέλιξη· τα παρακάτω είναι η
% κατάσταση 2026-07-18. Πριν την τελική εκδοχή του κεφαλαίου απαιτείται
% απόφαση χρήστη για το ποια ευρήματα γράφονται ως ΔΕΚΤΑ (validity review εκκρεμεί).

Για το φορτίο, η αξιολόγηση οργανώθηκε ως ``contest'' πέντε οικογενειών μοντέλων
(LGBM, XGBoost, MLP, LEAR, LSTM) έναντι της επίσημης day-ahead πρόβλεψης του
διαχειριστή (ΑΔΜΗΕ/ENTSO-E), σε τρία παράθυρα (Καλοκαίρι 2025, Οκτ.--Νοε. 2025,
Q1) και δύο πύλες πληροφορίας (g12: 12:00 CET D-1 gap-12h· g14: ευθυγραμμισμένη
με τον χρόνο δημοσίευσης της πρόβλεψης του διαχειριστή). Κύρια σημεία:

\begin{itemize}
  \item \textbf{Οκτ.--Νοε. 2025: το μοντέλο υπερτερεί του διαχειριστή.} Το
        XGBoost με καθαρά χαρακτηριστικά (χωρίς μετεωρολογικά, άρα χωρίς
        oracle πληροφορία) επιτυγχάνει MAE $137{,}5$--$139{,}6$~MW (g12) και
        $142{,}7$--$143{,}8$~MW (g14) σε τρεις seeds, έναντι $146{,}81$~MW του
        διαχειριστή. %TRACE: runs/load_contest{,_seeds}/octnov_xgb_recw_*dense*
        Το εύρημα είναι σταθερό σε 3 seeds και 2 πύλες, αλλά παραμένει
        window-specific: στο Q1 και το καλοκαίρι ο διαχειριστής διατηρεί
        σαφές προβάδισμα.
  \item \textbf{Πυκνές υστερήσεις φορτίου: θετικές σε όλους τους αλγορίθμους} ---
        LGBM, XGB, MLP, LEAR, σε 3 παράθυρα και 2 πύλες (XGB: 18/18 κελιά-seeds).
        %TRACE: ABLATION_PLAN §7.19/§7.21/§7.23
  \item \textbf{Έντιμος D-1 καιρός (meteo\_vintage).} Υλοποιήθηκε λήψη
        αρχειοθετημένων D-1 μετεωρολογικών προβλέψεων (Open-Meteo Previous Runs)
        με gate-aware ανάμειξη ημερών. Διατηρεί το $36$--$78\%$ του οφέλους του
        oracle καιρού --- δηλαδή τα συμπεράσματα που βασίζονταν σε oracle
        μετεωρολογικά υπερεκτιμούσαν το όφελος έως $\sim\!2{,}8\times$.
  \item \textbf{Στρατηγική: αλληλεπίδραση στρατηγικής--εποχής.} Σε πλήρη σύγκριση
        3 παραθύρων (LGBM, 5 σύνολα χαρακτηριστικών), η recursive υπερτερεί σε
        Q1 και Οκτ.--Νοε. (9/10 στατιστικά διακρίσιμα κελιά), ενώ η direct
        υπερτερεί το καλοκαίρι (4/5) --- η επιλογή στρατηγικής για το φορτίο
        είναι εποχιακά εξαρτημένη. Μόνο με την πρόβλεψη φορτίου του διαχειριστή
        στο σύνολο χαρακτηριστικών η recursive υπερτερεί και στα 3 παράθυρα
        ($-81{,}8$/$-84{,}1$/$-23{,}8$~MW). %TRACE: runs/load_direct/ §7.23-24, validity review Ε4
  \item \textbf{LSTM.} Μετά από διόρθωση θεμελιώδους σφάλματος exposure bias στο
        training loop, το LSTM με μόνο ημερολογιακά $+$ day-ahead πρόβλεψη
        φορτίου αναδείχθηκε ως η βέλτιστη διαμόρφωσή του (3/3 παράθυρα, 2 πύλες),
        με τα μετεωρολογικά και τις προβλέψεις ΑΠΕ να είναι επιζήμια στη
        συγκεκριμένη αρχιτεκτονική.
\end{itemize}

\section{Σύνοψη}
\label{sec:lf_summary}

Η leak-free εκστρατεία θεμελιώνει: (α) ένα επαληθεύσιμα έγκυρο κεντρικό
αποτέλεσμα για την τιμή DAM (LGBM \texttt{default,dense}, weekly, recursive ---
$16{,}96$~€/MWh στο Q1, $13{,}74$~€/MWh το καλοκαίρι, 3 seeds), (β) μια πρώτη
τεκμηριωμένη υπέρβαση της επίσημης πρόβλεψης φορτίου του διαχειριστή σε ένα
φθινοπωρινό παράθυρο με αυστηρά νόμιμη πληροφόρηση, και (γ) ένα επαναχρησιμοποιήσιμο
πρωτόκολλο εγκυρότητας (gate ανά πηγή, poisoning tests, κριτήρια αποδοχής
πολλαπλών παραθύρων/seeds) του οποίου η αξία καταδείχθηκε εμπειρικά από τις
διαρροές που εντόπισε και απέσυρε.
